% =========================================================================
% Capítulo 1 — Introducción
% Según la guía del TFG, la introducción debe exponer:
%   - los objetivos del TFG (ámbito, alcance y límites; aquí solo se
%     presentan: se desarrollan en su propio capítulo),
%   - el estado de la cuestión y las circunstancias de la elección,
%   - la pertinencia del tema y los motivos de su elección,
%   - la metodología empleada,
%   - la enumeración y descripción general de los capítulos, con su
%     justificación desde el punto de vista de la estructura.
% =========================================================================

\chapter{Introducción}
\label{ch:introduccion}

\section{El problema y la propuesta}
\label{sec:intro-problema}

Casi todo el mundo mantiene una lista, escrita o mental, de cosas que quiere hacer en la vida: viajes, hábitos, proyectos, experiencias —las llamadas \eng{bucket lists} o listas de deseos vitales—.
La mayoría de esas metas se queda en deseo: la intención por sí sola explica solo una parte pequeña de la conducta \parencite{sheeran-intention}, y las condiciones que la convierten en acción —plazos concretos, seguimiento del progreso, apoyo de otras personas— rara vez acompañan a una lista guardada en un cajón.
Las redes sociales dominantes no ayudan: su experiencia está volcada en el consumo de contenido, no en la acción —una lectura que el capítulo \ref{ch:beneficios} desarrolla y fundamenta—, y su uso pasivo está asociado a peor bienestar \parencite{verduyn-sns}.

\bukit es la respuesta que este \ac{TFG} construye: una red social cuyo propósito es ayudar a cumplir metas vitales y a no perseguirlas en solitario.
Toda la experiencia se orienta a la acción sobre tres pilares: estructura (metas agrupadas en listas con plazo y recordatorios), comunidad (personas con objetivos afines que se acompañan) y recompensa al avance (visibilidad para quien progresa, no para quien más audiencia acumula).
El capítulo \ref{ch:contexto} desarrolla el concepto completo y el \ref{ch:beneficios} fundamenta con literatura por qué ese diseño puede funcionar.

El resultado material del trabajo es una aplicación real: una \ac{API} propia con Node y PostgreSQL y un cliente web con mapa social, mensajería y verificación de identidad —el compromiso mínimo del trabajo—, acompañados como objetivos secundarios de una aplicación móvil parcial y de funciones asistidas por \ac{IA}, con la lógica compartida entre clientes desde paquetes comunes.
Los objetivos comprometidos, con su alcance y sus límites exactos, se definen en el capítulo \ref{ch:objetivos}.

\section{Pertinencia del tema y motivos de la elección}
\label{sec:intro-motivos}

Los motivos de la elección son tres.
El primero es la continuidad: el autor ya construyó y documentó TrackBites \parencite{deleanu-trackbites}, una aplicación web completa de registro nutricional, y este trabajo da el salto de dificultad que aquella experiencia deja preparado —de una aplicación web en JavaScript a un producto multiplataforma en TypeScript con base de datos relacional, geolocalización, búsqueda por afinidad y despliegue continuo (capítulo \ref{ch:alternativas})—.
El segundo es la cobertura de las competencias del grado: un producto de este tipo obliga a ejercitar en un solo proyecto el diseño de datos, la construcción de una \ac{API} segura, dos clientes con lógica compartida, las pruebas automatizadas, la entrega continua y la gestión completa del proyecto (planificación, riesgos, seguimiento y presupuesto) que esta memoria documenta.
El tercero es personal y el capítulo \ref{ch:beneficios} lo declara abiertamente: el autor quiere una herramienta que empuje a las personas de consumir a crear, y ese propósito existía antes que este \ac{TFG}.

\section{Metodología}
\label{sec:intro-metodologia}

El proyecto se gestiona con un método híbrido de la familia Scrumban adaptado a una sola persona: iteraciones semanales con capacidad variable declarada, tablero con límite de trabajo en curso, historias de usuario estimadas en puntos y medición continua de la velocidad para recalibrar el plan en cada cierre (capítulo \ref{ch:tareas}).
El desarrollo se realiza con asistencia declarada de \ac{IA} generativa, con salvaguardas explícitas contra la dependencia acrítica (riesgo R5, capítulo \ref{ch:riesgos}); cada decisión queda registrada en un diario de desarrollo con su alternativa descartada y su porqué, y esa trazabilidad es la materia prima de esta memoria.
La gestión de riesgos es un proceso continuo con seguimiento semanal (capítulo \ref{ch:riesgos}) y el gasto real del trabajo se registra mes a mes (capítulo \ref{ch:presupuesto}).

\section{Estructura de la memoria}
\label{sec:intro-estructura}

La memoria sigue el índice oficial de la \ac{EIB} y su lógica interna va de la justificación (qué se hace y por qué) a la solución (cómo), de ahí a la gestión (cuándo y a qué coste) y cierra con la evidencia de lo logrado.
El bloque de análisis lo forman el contexto y el concepto del producto (capítulo \ref{ch:contexto}), los objetivos y el alcance (capítulo \ref{ch:objetivos}), los beneficios y su relación con los \ac{ODS} (capítulo \ref{ch:beneficios}), el estado del arte con los requisitos del sistema (capítulo \ref{ch:estado-del-arte}), el análisis de alternativas tecnológicas (capítulo \ref{ch:alternativas}) y el plan de riesgos (capítulo \ref{ch:riesgos}).
El bloque de solución describe la propuesta (capítulo \ref{ch:solucion}) y su diseño de alto nivel (capítulo \ref{ch:diseno}).
El bloque de gestión y ejecución recoge la metodología con la planificación y la estimación del esfuerzo (capítulo \ref{ch:tareas}), los cálculos y algoritmos del producto (capítulo \ref{ch:algoritmos}), el análisis de resultados (capítulo \ref{ch:resultados}), el cronograma completo del \ac{TFG} (capítulo \ref{ch:gantt-tfg}) y el presupuesto (capítulo \ref{ch:presupuesto}).
Cierran las conclusiones y el trabajo futuro (capítulo \ref{ch:conclusiones}).
